\begin{abstract}
Single-cell transcriptomics data are sparse and noisy, motivating metacell approaches that aggregate similar cells to recover underlying biological states. However, in multi-batch settings, batch effects make similarity relationships reliable primarily within individual datasets, limiting the ability to aggregate cells across batches, even though many biological states are shared across them. At the same time, incorporating signals from other modalities such as spatial context while remaining faithful to gene expression is a challenge that existing methods do not address jointly. We introduce scProto, a framework that learns metacells by preserving structure encoded in a cell--cell affinity graph encoding biologically meaningful similarity (e.g., density-corrected transcriptomic or spatial), while remaining faithful to gene expression data. scProto combines a reconstruction objective, which ensures that learned representations capture meaningful gene expression patterns, with a community-preserving objective that encourages cells connected in the affinity graph to map to shared prototypes. Experiments on single-cell and spatial transcriptomics datasets show that scProto improves cross-batch aggregation while preserving biological structure, achieving higher batch mixing while maintaining comparable modularity and cell-type purity. It also better captures rare cell populations and spatially driven niches, demonstrating the benefit of affinity-guided clustering.
\end{abstract}
